\section*{第 \ref{ch:b1:plant-transport} 章——植物的气体交换与汁液运输}

\begin{solution}{exo:b1:plant-transport:1}
面向\omterm{def:b1:plant-transport:stoma}{气孔}的壁较厚，微纤丝像箍一样环绕\omterm{def:b1:cell-unit-of-life:cell}{细胞}排列，
因此膨胀的\omterm{def:b1:cell-unit-of-life:cell}{细胞}不能变宽，只能伸长；
两个这样的\omterm{def:b1:cell-unit-of-life:cell}{细胞}在两端相连，像一对香肠一样向外弯曲，
两\omterm{def:b1:cell-unit-of-life:cell}{细胞}之间的间隙随之张开。
\end{solution}

\begin{solution}{exo:b1:plant-transport:2}
叶片湿润的\omterm{def:b1:eukaryotic-cell:plantcell}{细胞壁}中的水分蒸发到气腔中，
由此产生拉力：\omterm{def:b1:eukaryotic-cell:plantcell}{细胞壁}细孔中的表面张力使水处于张力之下。
水的内聚力——即其\omterm{def:b1:water-small-molecules:water}{氢键}——
将这种张力沿\omterm{def:b1:flowering-plant-organization:tissues}{木质部}中连续不断的水柱向下传递。
在根部，张力使\omterm{def:b1:flowering-plant-organization:tissues}{木质部}的\omterm{def:b1:membranes-transport:osmosis}{水势}低于土壤，
从而将水分吸入。
\end{solution}

\begin{solution}{exo:b1:plant-transport:3}
约正午时；约 13:00；因为中午干燥的空气使
蒸腾作用快于根系供水的速度，叶片的
\omterm{def:b1:membranes-transport:osmosis}{水势}下降，\omterm{prop:b1:plant-transport:opening}{脱落酸}和叶片自身的水分胁迫
使\omterm{def:b1:plant-transport:stoma}{气孔}部分关闭。
\end{solution}

\begin{solution}{exo:b1:plant-transport:4}
源：输出糖分的\omterm{def:b1:mammal-organization:organ}{器官}（如叶片、萌发的块茎），在这里
蔗糖被装载入\omterm{def:b1:plant-transport:phloem}{筛管}，水随之进入，\omterm{def:b1:eukaryotic-cell:plantcell}{膨压}升高。
库：输入糖分的\omterm{def:b1:mammal-organization:organ}{器官}（如根、果实、生长顶端、正在充实的块茎），
在这里蔗糖被卸出，水分离开，\omterm{def:b1:eukaryotic-cell:plantcell}{膨压}下降。
\end{solution}

\begin{solution}{exo:b1:plant-transport:5}
$1000\times 9.8\times 30 = \qty{0.29}{MPa}$；计入摩擦后
\qty{0.6}{MPa}：\omterm{def:b1:flowering-plant-organization:tissues}{木质部}处于 \qty{-0.6}{MPa}，叶片则低于此值，
约 \qty{-1}{MPa}。
\end{solution}

\begin{solution}{exo:b1:plant-transport:6}
截面积相同，即 100 个\omterm{def:b1:flowering-plant-organization:tissues}{管胞}。每条管道的流量 $\propto r^4$，
一个\omterm{def:b1:flowering-plant-organization:tissues}{导管}为 $10^4$ 单位，每个\omterm{def:b1:flowering-plant-organization:tissues}{管胞}为 1 单位，一百个\omterm{def:b1:flowering-plant-organization:tissues}{管胞}合计 100，
因此\omterm{def:b1:flowering-plant-organization:tissues}{导管}的输送能力是\omterm{def:b1:flowering-plant-organization:tissues}{管胞}的一百倍。\omterm{def:b1:flowering-plant-organization:tissues}{管胞}较不易发生\omterm{def:b1:plant-transport:xylem}{空穴化}，
其中一个出现气泡时只会使一条很小的通道失效；
针叶树全部为\omterm{def:b1:flowering-plant-organization:tissues}{管胞}，能够度过冰冻的冬季，
而冬季冰冻会使含\omterm{def:b1:flowering-plant-organization:tissues}{导管}树木的\omterm{def:b1:flowering-plant-organization:tissues}{木质部}大量空化。
\end{solution}

\begin{solution}{exo:b1:plant-transport:7}
$\Psi_p = \Psi_{\text{木质部}} - \Psi_s = -0.5 + 1.5 = \qty{1.0}{MPa}$
在源端；在库端理论上为 $-0.5 + 0.4 = \qty{-0.1}{MPa}$
（管道需要略带张力，即实际\omterm{def:b1:eukaryotic-cell:plantcell}{膨压}约为零）；
约 \qty{1}{MPa} 的差值驱动水流。
\end{solution}

\begin{solution}{exo:b1:plant-transport:8}
截面积 $\pi\times 10^{-10} = \qty{3.1e-10}{m^2}$；\qty{1}{\micro L}
$= \qty{1e-9}{m^3}$ 每小时，速度 $3.2\,\mathrm{m/h}$——
偏高，因为口针处的流动快于
未受干扰的管道。糖分：每小时 \qty{0.2}{mg}。
\end{solution}

\begin{solution}{exo:b1:plant-transport:9}
环剥切断了\omterm{def:b1:plant-transport:phloem}{韧皮部}，糖分因此无法到达根部，
根部在数月内因饥饿而死，随后吸水功能丧失，
树冠随之死亡。木材中的\omterm{def:b1:flowering-plant-organization:tissues}{木质部}未受损伤，水分仍能到达叶片，
叶片可继续进行数周的光合作用——
糖分在环剥处上方积累，使该处树皮肿胀。
\end{solution}

\begin{solution}{exo:b1:plant-transport:10}
梯度：水 $23 - 9 = \qty{14}{g/m^3}$，$\mathrm{CO_2}$ $0.72 -
0.45 = \qty{0.27}{g/m^3}$；按质量计，通量之比为 $1.6\times 14/0.27 = 83$；
按分子计为 $83\times 44/18 = 200$ 个水分子，
对应每个 $\mathrm{CO_2}$。
\end{solution}

\begin{solution}{exo:b1:plant-transport:11}
水柱断裂：气泡膨胀，充满\omterm{def:b1:flowering-plant-organization:tissues}{导管}，
其上方的水不再被拉动，\omterm{def:b1:flowering-plant-organization:tissues}{导管}停止输导；水流转向相邻\omterm{def:b1:flowering-plant-organization:tissues}{导管}，
使它们的张力和风险随之增加；
由其供水的叶片部分萎蔫。\omterm{def:b1:flowering-plant-organization:tissues}{导管}之间的纹孔过细，
气泡的弯月面无法通过，因此气泡被限制在一条\omterm{def:b1:flowering-plant-organization:tissues}{导管}内。
夜间张力消失，气泡可能在\omterm{prop:b1:plant-water-minerals:rootpressure}{根压}作用下溶解；
春季，树木长出一圈新\omterm{def:b1:flowering-plant-organization:tissues}{导管}，
并弃用旧\omterm{def:b1:flowering-plant-organization:tissues}{导管}。
\end{solution}

\begin{solution}{exo:b1:plant-transport:12}
\omterm{def:b1:flowering-plant-organization:tissues}{木质部}自下而上运输水分和无机盐，其驱动力来自叶片处的蒸腾作用——
即太阳提供的能量，而非植物自身的能量；
\omterm{def:b1:plant-transport:phloem}{韧皮部}将糖分（及信号）从源运往库，
其驱动力是植物在两端通过消耗 \omterm{def:b1:water-small-molecules:atp}{ATP} 装载与卸出蔗糖
所造成的压力差，两端之间的流动无需耗能。
植物只在界面处耗能：\omterm{def:b1:plant-transport:stoma}{保卫细胞}的\omterm{prop:b1:plant-water-minerals:uptake}{质子泵}，
\omterm{def:b1:plant-transport:phloem}{韧皮部}装载器，以及为承受张力而构建的木材。
其间的一切——
均依靠物理原理运行。
\end{solution}

\begin{solution}{pb:b1:plant-transport:1}
\textbf{1.} $1000\times 9.8\times 100 = \qty{0.98}{MPa}$。
\textbf{2.} 约 \qty{2.0}{MPa} 的张力。
\textbf{3.} $\Psi = -0.1 - 2.0 = \qty{-2.1}{MPa}$（土壤加上两处
张力）。
\textbf{4.} 低于 \qty{-2.1}{} 时：取 $\Psi_s = -2.5$，\omterm{def:b1:eukaryotic-cell:plantcell}{膨压}至多为
$-2.1 + 2.5 = \qty{0.4}{MPa}$——一片整个
下午都接近萎蔫边缘的叶片。
\textbf{5.} 在 \qty{10}{m} 处：重力 \qty{0.1}{}，摩擦 \qty{0.1}{}，
\omterm{def:b1:flowering-plant-organization:tissues}{木质部} \qty{-0.3}{MPa}；叶片\omterm{def:b1:eukaryotic-cell:plantcell}{膨压}可达 $-0.3 + 1.5 = \qty{1.2}{MPa}$：
\omterm{def:b1:eukaryotic-cell:plantcell}{膨压}充足。
\textbf{6.} 在 \qty{1000}{m} 高度，树冠处张力将达 \qty{20}{MPa}，
叶片需要低于此值的 $\Psi_s$ 才能吸水：
对活\omterm{def:b1:cell-unit-of-life:cell}{细胞}而言不可能，任何\omterm{def:b1:flowering-plant-organization:tissues}{导管}都会在
第一个干燥的午后之前早已发生\omterm{def:b1:plant-transport:xylem}{空穴化}；极限由\omterm{def:b1:plant-transport:xylem}{空穴化}
和叶片决定，而非水的强度。
\textbf{7.} \qty{50}{L/h}；$\qty{1.4e-5}{m^3/s}$。
\textbf{8.} $0.2\times 0.5 = \qty{0.1}{m^2}$。
\textbf{9.} $1.4\times 10^{-5}/0.1 = \qty{1.4e-4}{m/s}$，即 \qty{0.5}{m/h}。
\textbf{10.} $0.1\times 100 = \qty{10}{m^3}$：树干中悬挂着十吨水。
\textbf{11.} $100/0.5 = \qty{200}{h}$，即八天。
\textbf{12.} $50\,000/2000 = \qty{25}{g}$ 每平方米每小时；
$25/18/3600 = \qty{0.39}{mmol}$ 每平方米每秒。
\textbf{13.} $(-2.1 + 0.1)/100 = \qty{-0.02}{MPa/m}$，是
重力梯度 \qty{0.01}{MPa/m} 的两倍。
\textbf{14.} $2000\times 4\times 10^{-6}\times 43\,200 = \qty{346}{mol}$。
\textbf{15.} $346/12 = \qty{28.8}{mol}$ 蔗糖，\qty{9.9}{kg}。
\textbf{16.} \qty{4.9}{kg} 蔗糖，存于 $4.9/0.15 = \qty{33}{kg}$ 的
汁液中。
\textbf{17.} \qty{31}{L}；以 \qty{0.6}{m/h} 计，\qty{24}{h} 内的流量
为 $\qty{3.6e-7}{m^3/s}$，对应速度 $\qty{1.7e-4}{m/s}$：截面积
\qty{2.2e-3}{m^2}，即 \qty{22}{cm^2} 的\omterm{def:b1:plant-transport:phloem}{筛管}——相当于一个薄的
树皮圆柱体。
\textbf{18.} $100/0.6 = \qty{170}{h}$，约一周。
\textbf{19.} $600\,000/18 = \qty{33000}{mol}$ 的水，用于
\qty{346}{mol} 的 $\mathrm{CO_2}$：约为 96 比 1 ——对 C3 植物而言偏低，
因为红杉的树冠生活在湿润的沿海空气中。
\textbf{20.} $-1.2 + 1.0 = \qty{-0.2}{MPa}$，与\omterm{def:b1:flowering-plant-organization:tissues}{木质部}的
\qty{-2.1}{} 相比：水会从\omterm{def:b1:plant-transport:phloem}{韧皮部}流向\omterm{def:b1:flowering-plant-organization:tissues}{木质部}，因此\omterm{def:b1:plant-transport:phloem}{韧皮部}
无法在这些数值下从树冠\omterm{def:b1:flowering-plant-organization:tissues}{木质部}吸水；高大树木树冠处的装载
必须进一步降低\omterm{def:b1:plant-transport:phloem}{韧皮部}的 $\Psi_s$
（使汁液浓度更高），或者树冠\omterm{def:b1:flowering-plant-organization:tissues}{木质部}的张力小于
这一简单估算值——这是高大树木生理学中的一个真实难题。
\textbf{21.} $-0.3 + 0.4 = \qty{0.1}{MPa}$；差值 $1.0 - 0.1 =
\qty{0.9}{MPa}$，跨越 \qty{100}{m}。
\textbf{22.} 摩擦减半，降至 \qty{0.5}{}：张力为 \qty{1.5}{MPa}，
树冠\omterm{def:b1:membranes-transport:osmosis}{水势}为 \qty{-1.6}{MPa}；叶片重新获得\omterm{def:b1:eukaryotic-cell:plantcell}{膨压}
（\qty{0.9}{MPa}），\omterm{def:b1:plant-transport:xylem}{空穴化}风险降低，但碳固定
在\omterm{def:b1:plant-transport:stoma}{气孔}关闭的数小时内减半。
\textbf{23.} 树冠 $-1.0 - 2.0 = \qty{-3.0}{MPa}$：低于叶片的
$\Psi_s$，因此不可能有\omterm{def:b1:eukaryotic-cell:plantcell}{膨压}，树冠萎蔫；树木必须
关闭\omterm{def:b1:plant-transport:stoma}{气孔}（使树冠\omterm{def:b1:membranes-transport:osmosis}{水势}升至约 \qty{-2.5}{}，\omterm{def:b1:eukaryotic-cell:plantcell}{膨压}
为零），在树冠顶部落叶，并等待降雨。
\textbf{24.} 它们的\omterm{def:b1:membranes-transport:osmosis}{水势}是全树最低的：为保持
\omterm{def:b1:eukaryotic-cell:plantcell}{膨压}，它们积累更多溶质，\omterm{def:b1:eukaryotic-cell:plantcell}{细胞壁}也更厚；为限制
蒸腾，叶片更小；它们生长缓慢，因为其\omterm{def:b1:plant-transport:stoma}{气孔}
提早关闭。
\textbf{25.} 树冠处约有 \qty{2}{MPa} 的张力（\omterm{def:b1:membranes-transport:osmosis}{水势}
接近 \qty{-2.1}{MPa}），汁液以约 \qty{0.5}{m/h} 的速度
在边材中流动。
\end{solution}
